\begin{lemma}[Scheide exits are tree exits]
Every Scheide exit vertex $v$ is a tree exit in the broader sense of
\noderef{TreeExit}: for the missing color $\gamma$ and defective color
$\delta$ in its witness, the terminal segment from the outside end of the
full $(\delta,\gamma)$-chain to $v$ is a $T$-exit path.
\end{lemma}
\begin{proof}
Unpack the witness from \noderef{ScheideExitVertex}.  It supplies the full
simple incidence path $P$ from a tree vertex $u$ to an outside vertex $z$,
and identifies $v$ as the last tree vertex on $P$.  Take the subpath from
$v$ to $z$ and reverse it.  It is still simple, begins at the outside
vertex-node $z$, ends at the tree vertex-node $v$, and has positive
multigraph length.  By the last-tree-vertex clause, its only vertex-node
in $V(T)$ is $v$.  The outside endpoint of a two-color path is incident
with only the last one of the alternating colors, so it misses the other.
Consequently the reversed subpath satisfies every clause of
\noderef{TreeExitPath}, and hence witnesses
\noderef{TreeExit} for the same two colors.  The defective-color clause is
carried unchanged from the Scheide witness.
\end{proof}
